\section{寺院、两京与道路：宗教和城市怎样互相放大}
\label{sec:religion-urban-exchange}

长安和洛阳的城门并不只让诏书、官员和军队进出。译场、寺院、施主、写经者和远来僧侣也依赖同样的道路、坊市与官府许可。635年阿罗本获准译经、664年玄奘去世而弟子与译场仍在，说明都城可以把跨区域的文本和人变成政治上可见的活动；这不等于唐廷对所有信仰一概开放，更不能由一块碑或一部僧传推断城市居民共有的信念。城市的容纳力来自官署、市场和赞助网络的重叠，也受这些网络随时可能改变的限制。%
\footnote{阿罗本与玄奘见 ST-635-NESTORIAN、ST-664-XUANZANG-DEATH，依据《旧唐书》相关记载、景教碑、《大慈恩寺三藏法师传》与《续高僧传》。碑文和高僧传各有自我说明、典范化与成书时间的问题。}

武周把这类可见性推入继承政治。690年改唐为周时，受命请愿、符瑞和礼仪不是宫廷背后的装饰，而是把宗室、官僚、寺院资源和都城观众编入合法性语言的工具。705年唐号恢复后，官员任用、地方关系和宗教网络并没有随一个国号一同复位。由此看，宗教并非只能归入“思想史”：它提供了声望、文本和集会的形式，使政治语言能越过宫门；国家同样试图通过许可、礼制和人事把这些形式接回可管理的秩序。%
\footnote{建周与复唐见 ST-690-ZHOU-FOUNDATION、ST-705-TANG-RESTORATION。《旧唐书》《新唐书》则天皇后纪、中宗纪及《资治通鉴》可互校基本年序；符瑞的自发性、请愿者意愿和政变分工均受后出的政治叙事塑造。}

安史以后，城市与寺院的关系又被财政和战争重新估价。819年迎法门寺佛骨入宫，皇帝寻求的并非仅是私人敬信：舍利传统、宫廷礼仪、僧俗施主和长安的围观空间，使一件宗教物能够成为国家声望的事件。韩愈的《论佛骨表》保存了官僚反对的声音，却不是全社会意见调查。正因为寺院同时连接土地、捐施、铜料、旅人和识字网络，它既能进入皇室仪式，也会在财政紧张时成为政府想重新登记和调配的对象。%
\footnote{迎佛骨与韩愈被贬见 LT-819-037，依据《旧唐书·宪宗本纪下》、韩愈〈论佛骨表〉、《唐会要》卷47〈议释教上〉。佛骨真伪、仪式规模和民众反应不宜由文集和本纪作确定复原。}

845年的废寺因此不能只用“崇道”或“反佛”概括。诏令所指向的是寺院资产、僧尼身份、铜料和户籍，也是在军费与财政紧张中把流动人口和地方资源重新纳入国家分类的尝试；宣宗朝的恢复则显示它没有在各地留下同样持久的结果。河西的情形尤其能打断长安视角：张议潮据沙州后，寺院、军户与文书网络参与了地方秩序的重建，敦煌文书让这种联系偶尔显形，却不能替代整座城市或整个河西的社会账本。%
\footnote{会昌废佛、张议潮起事与归义军受节见 LT-845-049、LT-848-051、LT-851-052。可据《旧唐书·武宗本纪》《唐会要》卷47、张议潮传及敦煌文书；废寺、还俗和“十一州”的数字或称谓均须按地区和文种检验。}

后周955年的寺院整顿把同一难题带到五代的军政世界。朝廷希望取得铜料、劳力和可登记的人口，却无法从诏令本身知道每州怎样废并、谁失去生计、哪些施主和地方官作出变通。成都、杭州、金陵、福州和广州的寺院与市场则在不同区域政权下继续构成联系网络。国号的更替没有使这些城市停止交换；相反，港口、水路、寺院和地方精英的组合，解释了为什么区域秩序能在中原反复易代时维持自身的节奏。%
\footnote{后周整顿寺院见 LT-955-112，依据《旧五代史·周世宗纪二》《五代会要》卷26〈寺〉与《资治通鉴》。诏令能说明中央目标，不能量化地方落实、资产去向或个人信仰。}

\paragraph{城市网络的得失。}
两京与区域城市使宗教文本、人员和物资获得跨地域移动的条件，也让政权能够借礼仪和许可争取声望。它们的代价在于，寺院一旦同时成为资产、户籍和铜料的汇集处，便会承受国家财政的压力。研究者从正史、僧传、碑刻和敦煌文书中看到的都是不同位置的声音；把它们并读，能够重建关系的张力，不能补造一份全民宗教或城市收入的总账。
